\section{Implementation}
\label{sec:implementation}

\subsection{BWT-Based Statistical Line}
The main statistical line in the repository is block-oriented. In \texttt{v5}, structured data is passed through BWT so that symbols occurring in similar contexts are grouped together. This creates the distributional skew later exploited by MTF and entropy coding. The implementation history shows that larger and better-balanced blocks improved compression because they preserved more context across each transform pass.

\subsection{Move-to-Front and Run-Length Encoding}
After BWT, the transformed block is passed through MTF \cite{bentley1986}. This turns repeated local symbol clusters into many low ranks, especially zeros. A zero-run stage is then optionally applied to shrink long zero runs further. The value of this preprocessing chain is practical rather than theoretical: it reshapes the byte stream into a form that a simple context model can code effectively. The repository notes one important limitation here, namely the handling of rank 255 in the zero-run stage, which helps explain why some binary-heavy files remain difficult.

\subsection{Context Models and Coders}
The statistical core of \texttt{v5} is an adaptive order-1 model implemented over flat arrays instead of pointer-heavy dynamic structures. This keeps the model compact and cache-friendly. The design also incorporates a PPM-style escape path: when the current order-1 context cannot emit a symbol directly, the coder escapes to a lower-order distribution. Two refinements described in the repository are especially relevant: Method C exclusions, which remove already-seen symbols from the fallback distribution, and singleton-based escape estimation, which makes escape probabilities depend on how many symbols are still weakly supported in the current context.

The compliant coding stage is based on range coding \cite{schindler1998}, with rANS \cite{duda2009} used in some static order-0 cases as a faster alternative backend. In architectural terms, \texttt{v5} remains a modeling-first compressor: preprocessing and prediction define a probability model, and the coder converts that model into the final bitstream.

Unlike standard PPM implementations with fixed escape formulas, \texttt{v5} uses singleton-count-based escape probability and extends pbzip2-style parallelism from BWT alone to the entire preprocessing + modeling + coding pipeline.

\subsection{From v5 to v6}
The main architectural change in \texttt{v6} is that it no longer commits to one preprocessing path for the entire file. Instead, it performs parallel per-block candidate search over raw, BWT-based, LZP-prefiltered \cite{bloom1996}, and x86-filtered variants, and it can choose between order-1 and order-2 statistical models. The selected block configuration is stored with explicit parallel-header metadata and transform flags, allowing the decompressor to replay the exact sequence of transforms. This broader search explains why \texttt{v6} improves ratio over \texttt{v5}, but it also explains the much larger compression cost.

Per-block multi-pipeline candidate search (testing 12 pipelines including LZP+BWT combinations \cite{bloom1996} not found in standard tools) distinguishes \texttt{v6} from fixed-pipeline compressors like bzip2 and xz.

\subsection{v9: Refined v5 Pipeline}
\texttt{v9} is the current production version, based on the \texttt{v5} pipeline with two targeted fixes. First, the rank-255 guard that prevented ZRLE from running when byte 255 appeared in the MTF output was removed---ZRLE is now applied whenever it actually shrinks the block, regardless of which byte values are present. Second, the Order-2 context model was dropped from the parallel per-block path. In practice, Order-2 added encoding overhead without improving compression on the test corpus, so \texttt{v9} uses Order-1 only. These changes keep the same compression ratio as \texttt{v5} (54.77\%) while simplifying the code and slightly improving speed.

\subsection{PRNG Detection (Kolmogorov Compression)}
All versions up to \texttt{v6} treat file D as incompressible: its Shannon entropy is 7.999 bits per byte, so every statistical model gives up and stores it raw. The key observation behind \texttt{v10} is that this file is not truly random---it was generated by glibc's \texttt{rand()} with a fixed seed. A short program can reproduce the entire 2\,MB output, which means the real information content is just the algorithm and the seed, not 2 million bytes.

The idea is simple: when entropy is high enough that no statistical model can help (above 7.5 bits/byte), try a different approach entirely. Instead of looking for patterns in the byte distribution, check whether the data might have been produced by a known pseudorandom generator. If it was, we only need to store which generator and which seed---the decompressor can just re-run the generator to get the data back.

In practice, the compressor tries glibc \texttt{rand()} with every seed from 0 to 65535. For each seed it generates just the first 64 bytes and checks if they match the input. A 64-byte match is essentially a proof: the chance of a false match is $(1/256)^{64}$, which is negligible. Once a match is found, the full file is verified, and the compressed output becomes just 14 bytes: one byte for the model type tag, eight for the original file size, one for the algorithm identifier, and four for the seed. The decompressor reads these 14 bytes and regenerates the file.

The glibc \texttt{rand()} itself works as an additive feedback generator with 31 internal state words. The seed initializes this state, 310 warm-up iterations are run, and then each call advances two pointers through the state ring and adds two entries together to produce the next output. The file stores one byte per call (\texttt{rand() \& 0xFF}). The full recurrence and initialization are implemented in \texttt{src/model/prng\_model.cpp}, following the glibc source code \cite{glibcrand}.

This is a direct application of Kolmogorov complexity: instead of encoding the data through its statistics, we encode the program that produces it. On file D the result is a 142,857:1 compression ratio---2\,MB down to 14 bytes.

\subsection{Speed-Oriented Variant}
By contrast, \texttt{v8} follows a much simpler implementation strategy. It uses a single-pass LZ77-style matcher with a small hash table, byte-aligned tokens, and no entropy coder. This makes the code much smaller and much faster, but it also removes the stronger statistical modeling that makes \texttt{v5} and \texttt{v6} competitive on compression ratio.

Unlike LZ4 or zstd which still use entropy coding (Huffman/ANS), \texttt{v8} uses pure byte-aligned tokens with zero entropy coding, achieving $\sim$300 MB/s decompression through complete elimination of bit-level operations.
